\chapter*{Agradecimientos}

Hay personas sin las cuales este trabajo no habría sido posible, y hay otras sin las cuales el camino no habría tenido el mismo sentido. Este Trabajo de Fin de Grado es el resultado de muchas horas de esfuerzo, dudas, aprendizaje y constancia, pero también de todas las personas que, de una forma u otra, me han acompañado durante esta etapa. A todas ellas van dedicadas estas líneas.

A mi padre, por estar siempre presente. Por cada vez que me preguntó cómo iba el TFG, aunque muchas veces no entendiera del todo la respuesta, y aun así me escuchara con toda la atención del mundo. Gracias por sostenerme en los momentos en los que el cansancio pesaba más que las ganas, por confiar en mí incluso cuando yo dudaba, y por recordarme, con gestos sencillos, que era capaz de llegar hasta aquí. Mucho de lo que soy, de mi manera de esforzarme y de no rendirme, viene de ti.

A \textbf{Iana Miranda} y \textbf{Macarena Pereira}, por haberse convertido en amigas de verdad durante esta etapa. Gracias por las tardes de estudio, por las conversaciones que empezaban con apuntes y terminaban arreglando la vida, por las risas entre exámenes y por estar también cuando las cosas se hacían cuesta arriba. Haber compartido la carrera con vosotras ha hecho que todo fuese más llevadero, más bonito y mucho más mío. Me llevo no solo los recuerdos de estos años, sino la suerte de haber encontrado a dos amigas que han sabido acompañarme con cariño, paciencia y alegría. Ojalá sigáis estando cerca en todo lo que venga después.

A \textbf{Irene González} e \textbf{Ignacio Mora}, por sus aportaciones, su tiempo y su disposición a ayudar a mejorar este proyecto. Gracias por tomaros el interés que os tomasteis, por las ideas que aportasteis y por hacerme ver aspectos que yo sola no habría visto. Vuestra ayuda me permitió mirar el trabajo con más perspectiva y mejorarlo con más criterio.

A \textbf{Santia}, por haber compartido conmigo tantos momentos de esta etapa y por haber estado cerca también durante este último tramo. Gracias por las dudas compartidas, por los ánimos en los días complicados y por esas pequeñas victorias que se disfrutan más cuando alguien entiende lo que ha costado conseguirlas.

Y, de manera muy especial, a mi tutor, \textbf{José Antonio Parejo}, por su implicación, su cercanía y su paciencia durante todo el desarrollo de este trabajo. Gracias por aportar siempre ideas que ayudaban a mejorar el proyecto, por responder con rapidez cada duda y por guiarme con claridad cuando necesitaba ordenar el camino. Sus orientaciones han sido fundamentales para que este TFG avanzara con sentido y para convertir una idea inicial en un proyecto del que me siento orgullosa. Este trabajo también es suyo.
